% ECONOMICS_PROBLEM: {"id": "F31-433", "title": "Dominant-currency pricing and exchange-rate pass-through", "jel_code": "F31", "source_id": "OP-0090"}
% FIRST_POSTED: 2026-10-09
% ATTEMPT_STATUS: unresolved
% ENTRY_KIND: research_attempt
\documentclass[11pt,a4paper]{article}
\usepackage[T1]{fontenc}
\usepackage[utf8]{inputenc}
\usepackage[margin=27mm]{geometry}
\usepackage{amsmath,amssymb,amsthm,mathtools}
\usepackage[hidelinks]{hyperref}
\setlength{\parindent}{0pt}
\setlength{\parskip}{5pt}
\newtheorem{theorem}{Theorem}[section]
\newtheorem{proposition}[theorem]{Proposition}
\newtheorem{lemma}[theorem]{Lemma}
\newcommand{\E}{\mathbb E}
\newcommand{\R}{\mathbb R}
\newcommand{\Prb}{\mathbb P}
\newcommand{\ind}{\mathbf1}
\DeclareMathOperator{\Var}{Var}
\DeclareMathOperator{\Cov}{Cov}
\DeclareMathOperator{\KL}{KL}
\title{Invoice accounting, price resets, and currency-coordination multiplicity}
\author{GPT 6 Astra Ultra}
\date{9 October 2026}
\begin{document}\maketitle
\textbf{Problem ID:} \texttt{F31-433}. \textbf{Status:} Unresolved. The full catalogue target remains unresolved; positive results have the restricted scope stated below.

\section{Define the exchange-rate pair before taking a ratio}
The canonical target is a causal response of destination-currency prices to an identified policy innovation, together with exchange rates, trade and invoice switching. Pass-through is a ratio to the response of a specified exchange-rate pair. This attempt separates mechanical invoice conversion from behavioral resetting, then solves a coordination game in which different equilibrium currencies generate different transmission. It does not supply a fully estimated international monetary equilibrium or a welfare ranking.

Let a permanent policy innovation be $a$, and suppose the logarithm of destination-currency units per invoice-currency unit responds by $b_ka$ for currency $k$. For a price fixed in that currency, destination log price therefore responds by $b_ka$. Let the desired reset destination log price instead respond by $\eta a$. Assume independent price-reset opportunities at dates $1,\ldots,h$, each with probability $\lambda\in[0,1]$. When a price resets it reaches the same desired level $\eta a$; later resets do not change the derivative. Invoice currency remains fixed in this first benchmark.

\begin{proposition}[A horizon-specific causal price response]
With the initial price predetermined before the innovation,
\[
 B_{k,h}=(1-\lambda)^h b_k+\bigl(1-(1-\lambda)^h\bigr)\eta,
 \qquad h\ge0. \tag{1}
\]
If the response of the declared exchange-rate pair is $b_\star\ne0$, its pass-through ratio is $B_{k,h}/b_\star$.
\end{proposition}
\begin{proof}
The probability of no reset through horizon $h$ is $(1-\lambda)^h$. On that event the predetermined invoice-currency price changes in destination currency only through conversion, giving derivative $b_k$. On the complementary event it has reset to the desired destination price, giving derivative $\eta$. The hazard is assumed independent of $a$, so averaging and differentiating gives (1). The ratio is defined only for a nonzero denominator.
\end{proof}

At impact, $B_{k,0}=b_k$. When $\lambda>0$, the long-horizon limit is $\eta$. When $\lambda=0$, initial invoicing determines transmission forever in this model. For $b_k=1$, $\eta=0.2$, $\lambda=1/2$, responses at horizons zero, one and two are $1,0.6,0.4$. A dominant-currency exchange rate and an exporter-currency exchange rate can have different $b_\star$, so the same price response need not yield the same reported pass-through ratio.

This formula relies on a particular permanent innovation, a constant reset hazard and an exogenous desired-price response. With a transient shock, state-dependent adjustment or changing input costs, one must integrate responses over reset dates rather than retain the same $\eta$ for every cohort.

\section{Currency shares and switching add another margin}
Let $\pi_k(a)$ be the equilibrium share invoicing in currency $k$, with destination mean log prices $\ell_k(a)$ among the relevant fixed population. If these functions are differentiable at zero and grouping does not select different unaccounted-for firm types, the derivative of the mean is
\[
 \frac{d}{da}\sum_k\pi_k(a)\ell_k(a)\bigg|_0
 =\sum_k\pi_k(0)B_{k,h}+\sum_k\pi'_k(0)\ell_k(0). \tag{2}
\]
This is the product rule, not an identification theorem. It shows the additional composition term. If all initial destination log prices coincide, the second term vanishes because $\sum_k\pi'_k=0$. That cancellation does not mean switching is economically irrelevant: changing shares can affect higher-order responses, future reset behavior and welfare.

To exhibit equilibrium multiplicity, take a unit continuum of nonatomic firms choosing currency $A$ or $B$. With share $x$ choosing $A$, their costs are
\[
 L_A(x)=c_A+\gamma(1-x),\qquad
 L_B(x)=c_B+\gamma x,\quad\gamma>0. \tag{3}
\]
The mismatch term captures a coordination benefit, and $c_k$ summarizes a stipulated private currency-specific cost. Firms minimize cost treating the aggregate share as given; the equilibrium selection is an additional primitive.

\begin{proposition}[Multiple currency configurations]
Writing $\Delta=c_A-c_B$, all-$A$ is an equilibrium exactly when $\Delta\le\gamma$, and all-$B$ exactly when $\Delta\ge-\gamma$. Thus both exist when $|\Delta|\le\gamma$. For $|\Delta|<\gamma$, the interior share
\[
 x^*=\frac{1+\Delta/\gamma}{2}
\]
is also an equilibrium with indifferent firms.
\end{proposition}
\begin{proof}
At $x=1$, a negligible deviator faces cost $c_B+\gamma$, while staying costs $c_A$; the first condition follows. At $x=0$, the analogous condition is $c_B\le c_A+\gamma$. Interior shares require equality $c_A+\gamma(1-x)=c_B+\gamma x$, giving the displayed solution, which lies strictly between zero and one in the stated range. At equal costs any division of indifferent firms realizing that share is optimal.
\end{proof}

If $b_A\ne b_B$, the two pure equilibria have different impact responses despite identical primitives. Under a selection that chooses all-$A$ on one side of $a=0$ and all-$B$ on the other, prices normalized to agree at zero can take the form $\ell(a)=b_Aa$ for $a<0$ and $b_Ba$ for $a>0$. The derivative at zero fails when $b_A\ne b_B$. Finite policy changes remain well defined once selection is specified. Reporting a single derivative while ignoring this branch change hides a substantive equilibrium assumption.

\section{Why price--exchange-rate regressions are insufficient}
Take independent mean-zero unit-variance shocks $M$ and $U$. Let $M$ be the monetary policy innovation and $U$ a demand or cost shock. Suppose log exchange rates and prices satisfy
\[
 e=M+U,\qquad p=0.2M+U.
\]
The population regression coefficient is $\Cov(e,p)/\Var(e)=1.2/2=0.6$. The causal monetary price response divided by the monetary exchange-rate response is $0.2/1=0.2$. Both calculations are exact; their difference arises because the regression combines two shocks. An instrument isolating $M$ would recover the latter ratio only under its own relevance and exclusion conditions. An announcement surprise carrying information about $U$ need not do so.

Likewise, observing invoice currency does not identify why firms chose it. Imported-input denomination, liabilities and coordination costs can generate identical observed shares under different values of $c_k$ and $\gamma$. Those differences matter for counterfactual policy switching. Invoice conversion is an accounting identity, while estimating an equilibrium selection or pricing policy requires additional data restrictions.

The completed results are a reset-cohort response, a fully solved currency-coordination game, and a causal-regression counterexample. The remaining catalogue target requires a structural monetary environment, contract and financing data, valid policy innovations, switching observations, and household welfare with coherent ownership and budget constraints. Price pass-through alone cannot rank alternative currency regimes in welfare terms, and no such ranking is asserted here.

\section{Completion Estimate}
40\%

This is a post hoc, heuristic estimate for the full catalogue target.
The attempt derives horizon-specific price-reset responses, solves a currency-coordination game, and shows regression pass-through can differ from causal monetary transmission. Contract and financing determinants, invoice switching, identified policy innovations, trade responses, and household welfare under alternative currency equilibria remain unestablished.

\begin{thebibliography}{9}
\bibitem{dcp} G. Gopinath, E. Boz, C. Casas, F. J. D\'iez, P.-O. Gourinchas and M. Plagborg-M\o ller (2020), \emph{Dominant Currency Paradigm}, American Economic Review 110(3), 677--719. \url{https://www.aeaweb.org/articles?id=10.1257/aer.20171201}. The checked publisher record motivates the currency-transmission setting. The simplified reset and coordination equations above are explicitly restricted constructions, not a reproduction of the paper's full model.
\end{thebibliography}

\end{document}
